\section{CQ and its normal form}
\label{sec:cq}\label{sec:languages}

\paragraph{Grammar.}
CQ clauses use derived parsing; expressions, patterns and types are
\code{syn::Expr}, \code{syn::Pat} and \code{syn::Type}. Macros stay opaque
token trees~\cite{host-synmacro}.
\begin{lstlisting}
Program ::= struct Name; Decl* Rule+
Decl    ::= relation R(c: Type, ...) [= (Expr)];
Rule    ::= h(Expr, ...) :- Clause, ... ;
Clause  ::= R(Arg, ...) | Guard | let Pat [: Type] = Expr
          | for Pat in Expr | Block
Guard   ::= ["#[total]"] "if" Expr
Arg     ::= x | _ | Expr
Block   ::= { Body } ( | { Body } )+  |  { Body } ( else { Body } )+
\end{lstlisting}
Bare single-segment paths bind logical variables; other arguments compute.
Several rules for one head are sugar for one rule whose body is a top-level
union over fresh positional head columns; each rule gets fresh local names,
opaque syntax receives lexical aliases rather than token substitution.
A trailing \code{= (Expr)} supplies the host value.

\paragraph{Lexical dependencies.}
$\fv_S(e)=(F,o)$ returns visible value names and an opacity bit; with logical
environment $B$, $\uses(e,B)=B$ if $o$ else $F\cap B$, and the diagnostic
$\operatorname{known}(e,Q)=\emptyset$ if $o$ else $F\cap Q$. Write $e\checkmark_B$
for $\operatorname{known}(e,Q)\subseteq B$. Attributes, macros, local items and
verbatim syntax are opaque; qualified path segments and field/method names are
not value uses.
\begin{mathpar}
\Rule{FV-Name}{x\notin A}{A\vdash x\rightsquigarrow(\{x\},\mathsf{false})}

\Rule{FV-Bound}{x\in A}{A\vdash x\rightsquigarrow(\emptyset,\mathsf{false})}

\Rule{FV-Children}{A\vdash e_j\rightsquigarrow(F_j,o_j)\quad(1\le j\le n)}
 {A\vdash c(e_1,\ldots,e_n)\rightsquigarrow(\bigcup_jF_j,\bigvee_jo_j)}

\Rule{FV-Closure}{A\cup\binds(\bar p)\vdash e\rightsquigarrow(F,o)}
 {A\vdash |\bar p|e\rightsquigarrow(F,o)}

\Rule{FV-For}{A\vdash e\rightsquigarrow u\quad A\cup\binds(p)\vdash b\rightsquigarrow v}
 {A\vdash\mathsf{for}\ p\ \mathsf{in}\ e\ b\rightsquigarrow u\cup v}

\Rule{FV-CondLet}{A\vdash e\rightsquigarrow u}
 {A\vdash\mathsf{let}\ p=e\rightsquigarrow_c(u,A\cup\binds(p))}

\Rule{FV-And}{A\vdash e_1\rightsquigarrow_c(u,A_1)\quad A_1\vdash e_2\rightsquigarrow_c(v,A_2)}
 {A\vdash e_1\mathbin{\&\&}e_2\rightsquigarrow_c(u\cup v,A_2)}

\Rule{FV-Cond}{A\vdash e\rightsquigarrow u\quad e\text{ is neither case above}}
 {A\vdash e\rightsquigarrow_c(u,A)}
\end{mathpar}
Blocks visit initializers and let-else branches in the old scope, then add
pattern bindings; match arms add their pattern bindings; condition bindings
reach only the then/loop body.
\begin{lemma}[Lexical upper bound]
For transparent syntax with the intended pattern resolution, every free simple
value name outside $A$ is in $F$, so $\uses(e,B)$ contains the logical free names
of $e$ in $B$; for opaque syntax this holds because all of $B$ is returned.
\end{lemma}
\begin{proof}Simultaneous induction on expressions, statements and conditions;
each binder extends only its body scope. Macro expansion and pattern resolution
remain hypotheses.\end{proof}

\paragraph{Scope rules.}
$\Sigma;B\vdash c\Rightarrow B'$ checks one clause in written scope; atom
arguments fold left to right.
\begin{mathpar}
\Rule{CQ-Var}{}{\Sigma;B\vdash x\Rightarrow B\cup\{x\}}

\Rule{CQ-Expr}{e\checkmark_B}{\Sigma;B\vdash e\Rightarrow B}

\Rule{CQ-Atom}{R\in\dom\Sigma\quad |\bar a|=|\Sigma(R)|\quad
 \Sigma;B\vdash\bar a\Rightarrow B'}{\Sigma;B\vdash R(\bar a)\Rightarrow B'}

\Rule{CQ-Guard}{e\checkmark_B}{\Sigma;B\vdash\mathsf{if}\ e\Rightarrow B}

\Rule{CQ-Bind}{e\checkmark_B\quad N=\set(\binds(p))\quad N\cap B=\emptyset}
 {\Sigma;B\vdash\mathsf{let}\ p[:\tau]=e\Rightarrow B\cup N}

\Rule{CQ-For}{e\checkmark_B\quad N=\set(\binds(p))\quad N\cap B=\emptyset}
 {\Sigma;B\vdash\mathsf{for}\ p\ \mathsf{in}\ e\Rightarrow B\cup N}

\Rule{CQ-Block}{n\ge2\quad\forall j,\ \Sigma;B\vdash B_j\Rightarrow B_j'\\
 I=(\bigcap_j B_j')\setminus B\quad \forall x\in I,\ \tau_j(x)=\tau_k(x)}
 {\Sigma;B\vdash\mathsf{block}(B_1,\ldots,B_n)\Rightarrow B\cup I}
\end{mathpar}
The generator RHS is checked in the old scope: its own pattern names are not
in scope until subsequent clauses (and the head). As for scalar binds, rustc
checks an irrefutable pattern and concrete binding types; for generators the
pattern destructures \code{IntoIterator::Item}, not the collection.
There is no surface generator type annotation. Finiteness and purity are
semantic premises, not properties established by the structural scope check.
An atom using a generator-supplied name requires that name already bound.


\paragraph{Recognized comparisons.}
Let $Q$ be the query binders, $B$ the written bound names and $P$ the closing
host parameters. Over the Syn AST (no separate grammar):
\begin{lstlisting}
t ::= x | *x | &x | literal | -numeric-literal | bare-host-identifier
C ::= t == t | t != t | t < t | t <= t | t > t | t >= t
    | C && C | C || C | !C | (C)
\end{lstlisting}
Rust desugars these to \code{PartialEq}/\code{PartialOrd} calls; we recognize
them only under the premises.
\begin{mathpar}
\Rule{ORD-Admissible}{\tau:\code{Ord}\quad
 \leq_\tau\text{ total, terminating and consistent with identity}}
 {\tau\text{ admits relational order}}

\Rule{CQ-Comparison}{\vars(t_1,t_2)\subseteq B\quad
 t_1,t_2\text{ admissible}\quad op\in\{==,\ne,<,\leq,>,\geq\}}
 {\Sigma;B;P\vdash t_1\ op\ t_2\Rightarrow B}

\Rule{CQ-GuardRecognize}{e\text{ has only recognized comparisons and boolean nodes}\\
 \text{all comparisons satisfy the Eq/Ord premises}}
 {\mathsf{recognize}_{Q,P}(e)=\theta\quad\sem e=\sem\theta}

\Rule{CQ-GuardSplit}{e=e_1\mathbin{\&\&}\cdots\mathbin{\&\&}e_n\\
 A=\text{recognized leaves}\quad F=\text{remaining leaves in written order}}
 {\mathsf{if}\ e\rightsquigarrow \mathsf{select}_{\bigwedge A};
   \mathsf{guard}_{\bigwedge F}\quad
   \sem{\mathsf{select}_{\bigwedge A};\mathsf{guard}_{\bigwedge F}}
   \sqsupseteq\sem{\mathsf{if}\ e}}
\end{mathpar}
Recognized comparisons are total, so short-circuit \code{\&\&}, \code{||} and
\code{!} mean classical connectives on them. Except for CQ-Total below,
arithmetic, calls, methods, casts, indexing, blocks, macros, qualified paths
and field access are opaque; a mixed \code{||} or \code{!} is one opaque leaf;
only the top-level conjunction splits.
Parameters print as \code{@name} in later languages and are borrowed once outside
query loops.
\begin{proof}[Proof of CQ-GuardSplit]
In \code{a \&\& f(x)} Rust evaluates $f$ exactly where $a$ holds, so selecting
the total $a$ earlier gives that domain; in \code{f(x) \&\& a} it evaluates $f$
on a superset, so selecting $a$ earlier only shrinks it. Iterating over the
conjunction keeps the opaque leaves in written order: on a returning input,
either a relational leaf rejects first or all hold and the residual evaluates
the same opaque leaves in the same order. A disjunction or negation lacks this
argument and is never split.
\end{proof}

\paragraph{Total Rust predicates.}
\[
\Rule{CQ-Total}{e\text{ is allowlisted over declared primitive types, or explicitly asserted total}\\
 F=\uses(e,B)}
 {[\code{\#[total]}]\ \mathsf{if}\ e\rightsquigarrow
  \mathsf{constraint}\ \code{\#[total]}\{x=x\mid x\in F\}(e)}
\]
The allowlist admits standard \code{str}/\code{String}
\code{contains}, \code{starts_with}, and \code{ends_with} with literal or
declared string arguments, primitive comparisons, Boolean literals, and
\code{\&\&}/\code{||} compositions. Unknown column types, calls, arithmetic,
mutable references, attributes and macros are not inferred total.
Transparent macro-forwarding groups preserve the declared-type proof, and
lexical predicate bindings retain the original identifier hygiene.
\code{\#[total] if e} is a trusted user premise that $e$ cannot diverge or panic,
in addition to the existing pure functional Rust premise. Hidden macro uses
conservatively retain the complete logical environment.
The normalized predicate retains the original Rust expression and explicit
lexical-to-relational bindings; aliases change bindings, not Rust tokens.
Its totality permits ordinary selection motion; unrecognized guards still
belong to the exact Rust chain. No data statistics or key assumptions enter
this classification.
The marker is an outer guard attribute, not a reserved identifier:
\code{total} remains available as a query variable, let binding and Boolean
guard operand. Normalized predicates also use the attribute-prefixed wire
form, so a relational comparison such as \code{total > 0} is unambiguous.

\paragraph{Normal form.}
CQ$_{\rm NF}$ restricts the same \code{cq::Program}: every atom argument is a
distinct bare variable, every head argument a variable. Shared names across
atoms are natural-join variables. Repeats within an occurrence and literals
become relational identities at the atom's position; computed arguments become
positioned binds plus identities; computed heads become final binds.

\paragraph{Contract.}
$\mathsf{CQScope}(s)$ is the scope rules plus:
(1)~relation and column names unique, declared arity positive, no attributes;
(2)~body atoms name declared inputs at their arity;
(3)~lets and generators do not shadow bound names; the body has an atom; the head has positive arity;
(4)~a declaration named by the head is the result: matching arity, no host value;
without it, head arguments are bound bare variables of known type;
(5)~blocks have equal interface types and no escaped local names;
(6)~negation, aggregation and unexpanded \code{name!\{..\}} clauses are rejected;
(7)~comparison terms are bound names, literals or captured identifiers; known
Eq/Ord failures are rejected; a recognizable unsplit guard is rejected.
\begin{mathpar}
\Rule{WF-CQ}{\mathsf{CQScope}(s)}{\mathsf{WF}_{CQ}(s)\Downarrow\ok}

\Rule{WF-NF}{\mathsf{WF}_{CQ}(s)\quad\text{every nested atom has distinct bare arguments}\\
 \text{head arguments are bound variables; known inadmissible types rejected}}
 {\mathsf{WF}_{NF}(s)\Downarrow\ok}
\end{mathpar}
These checks do not establish admissibility, which remains a premise.

\paragraph{Meaning and rewrite check.}
The meaning is Section~\ref{sec:denotation}: compose the desugared body in
written order. Its total relational specialization agrees with classical
$\sem q_I$. CQ records no decisions. A same-language CQ rewrite must be AST
identity; raw-token surface expansions run before the CQ contract and have no
semantic contract of their own.
